\chapter[Relaxed Branch, Export Kernel, and Material Thresholds]{Relaxed Open-System Branch, Barrier Compiler, Export Kernel, and Material Thresholds}
\label{chap:part08_relaxed_branch}

\section*{Stage range: 243--253}
\addcontentsline{toc}{section}{Stage range: 243--253}

Part VII ended at the strict selected-support and rigid-mouth front end.  The exact one-port same-charge audit had already fixed the admissible short-range kernel families, the actual-branch packet had already been compressed to orbit-lock variables, and the selected passive/outgoing support side had already landed on the lowest-twin support slice.  Part VIII is deliberately different.  It is not another attempt to prove the strict conservative branch.  It is the relaxed/open-system companion that asks what happens when three standard-recovery suppressions are lifted in a controlled way around the Stage 242 front end.

The conceptual firewall is the most important part of this chapter.  The relaxed branch does \emph{not} undo the same-charge verdict from Part VII.  It does not produce a new long-range same-charge attraction, and it does not introduce a new linear dynamic kernel class.  The relaxed branch is a codimension-three lift of the strict selected-support/orbit packet.  Its standard-recovery slice is
\begin{equation}
\label{eq:part08-intro-recovery-slice}
\ell_w=0,
\qquad
f_U=0,
\qquad
a=b=0,
\end{equation}
which implies
\begin{equation}
\label{eq:part08-intro-recovery-observables}
S_{\rm leak}=0,
\qquad
\mathcal W_w=0,
\qquad
U=0,
\qquad
V=0,
\qquad
\mathcal D_{UV}=0,
\qquad
\varsigma\equiv 1.
\end{equation}
Away from this slice, the relaxed branch can reshape a near-contact barrier by leakage, scalar-photon work, non-rigid dressing drain, and compensated source response.  Those are branch-local open-system or support/dressing/source effects.  They do not alter the long-range kernel firewall.

The central stationary object of this part is the lowered front end
\begin{equation}
\label{eq:part08-intro-veff230}
V_{\rm eff}^{(247)}(r)
=
V_{\rm short}^{(1p)}(r)
-
\lambda_L S_{\rm leak}(r)
-
\lambda_W\mathcal W_w^{\rm sess}(r)
-
\Delta E_{UV}(r)
-
\mathcal M_{\sigma}(r).
\end{equation}
The central dynamic object is the event-chain energy law
\begin{equation}
\label{eq:part08-intro-event-energy}
E=\frac{m_s}{2}\dot r^{\,2}+V_{\rm eff}^{(247)}(r).
\end{equation}
The central microscopic export object is the super-Ohmic active-leg kernel
\begin{equation}
\label{eq:part08-intro-export-kernel}
K_{\rm exp}(s)=\Gamma_3s^3+\Gamma_5s^5.
\end{equation}
The central materials companion is the four-ratio screen
\begin{equation}
\label{eq:part08-intro-material-screen}
\Pi_{\rm ep}\ge 1,
\qquad
\Pi_{\chi}\ge 1,
\qquad
\Pi_k\ge 1,
\qquad
\Pi_T(T)\ge 1.
\end{equation}

\section*{Part VIII at a glance}

\begin{readermap}
\item[Primary question] If the strict same-charge branch is kept fixed, what controlled open-system extensions can be built around it without undoing the strict kernel verdict?
\item[Starting inputs] The Part VII strict front-end packet, including the rigid-mouth orbit-lock variables, the selected twin-support slice, and the short-range kernel firewall.
\item[Delivers] The relaxed-constraint branch declaration, leakage/work and non-rigid \((U,V)\) compilers, the relaxed barrier and event-chain diagnostics, the microscopic export kernel, and the material-threshold companion.
\item[Used by] Future relaxed-branch, export-kernel, and material-threshold papers; no later part in the current ledger supersedes this handoff block.
\item[Result anchors] \resultanchor{MTDC-T12}.
\end{readermap}

\paragraph{Stage packets.}
\begin{center}
\small
\begin{tabular}{@{}>{\raggedright\arraybackslash}p{0.11\textwidth}>{\raggedright\arraybackslash}p{0.24\textwidth}>{\raggedright\arraybackslash}p{0.19\textwidth}>{\raggedright\arraybackslash}p{0.38\textwidth}@{}}
\toprule
Stages & Packet & Status & Carry-forward output \\
\midrule
243--246 & Relaxed branch declaration and lane compilers & \StatusEffective{} / \StatusExactClosure{} & Declares the codimension-three lift of the strict Stage 242 front end, fixes the exact recovery slice, and compiles the leakage/work, non-rigid mouth, and compensated multimode source lanes. \\
\addlinespace
247--250 & Lowered barrier, event-chain, and Goldilocks diagnostics & \StatusExactClosure{} / \StatusEffective{} / \StatusReduced{} / \StatusNumerical{} & Builds the relaxed stationary barrier, the one-dimensional event-chain and tunneling diagnostics, the helicity-export readback, and the crossing-versus-collapse window used to identify viable relaxed events. \\
\addlinespace
251--253 & Export kernel, heat partition, and material screening & \StatusExactClosure{} / \StatusReduced{} / \StatusEffective{} / \StatusNumerical{} / \StatusOpen{} & Derives the cubic/quintic active-leg export kernel, partitions exported energy between vacuum and lattice channels, and compiles the physical calibration and material-threshold screen without claiming that the completed PDE realizes those open-system coefficients. \\
\bottomrule
\end{tabular}
\end{center}

\claimstatus{Part VIII is a relaxed/open-system companion.  Its algebraic compilers are \StatusExactClosure{} inside the declared branch, one-mode, characteristic-timescale, single-growth, and physical-calibration closures.  The stationary and dynamic readbacks that use Session-I through Session-V benchmark numbers are \StatusNumerical{} insertions into exact formulas.  The branch declaration itself is an \StatusEffective{} extension around the strict Stage 242 front end, and any claim that the completed nonlinear moving-throat PDE actually realizes the relaxed packet, the microscopic export coefficients, or the physical material thresholds remains \StatusOpen{}.  No result in this part should be cited as a proof of a new same-charge long-range attraction or as a proof that the relaxed branch validates the strict conservative branch.}

\derivationfield{What this part proves}{Part VIII constructs the relaxed/open-system companion around the strict Stage 242 packet, including leakage/work, non-rigid mouth, compensated source, barrier, export, heat-partition, and material-screening compilers.}
\derivationfield{What this part does not prove}{It does not rescue the strict same-charge conservative branch, does not create a new long-range same-charge attraction, and does not prove that the completed PDE realizes the relaxed packet or its physical screening thresholds.}
\derivationfield{Why later parts need it}{This is the handoff block for future relaxed-branch, export-kernel, and material-threshold papers; no later part inside the current ledger supersedes it.}

\section{Block contract and theorem-level output}
\label{sec:part08-block-contract}

The contract of Part VIII is to keep three layers separate.  The first layer is the strict front end from Part VII: selected twin-support placement, rigid-mouth orbit lock, and the outgoing normalization packet.  The second layer is the relaxed branch around that front end: leakage, non-rigid mouth response, and compensated source shape.  The third layer is the physical companion: event-chain survival, export, heat partition, and material screening.  The chapter is useful only if those layers remain visibly separated.

The operational chain is
\begin{equation}
\label{eq:part08-operational-chain}
\begin{aligned}
\mathcal P_{242}
&\longrightarrow
\mathfrak B_{243}^{\rm relax}
\longrightarrow
(S_{\rm leak},\mathcal W_w,U,V,\mathcal M_\sigma)
\\
&\longrightarrow
V_{\rm eff}^{(247)}
\longrightarrow
(r_\pm,I_{\rm new},t_{\rm cross})
\\
&\longrightarrow
K_{\rm exp}(s)
\longrightarrow
(\Pi_{\rm ep},\Pi_\chi,\Pi_k,\Pi_T).
\end{aligned}
\end{equation}
The strict branch can be recovered from this chain by imposing \cref{eq:part08-intro-recovery-slice}.  The relaxed branch cannot be used to infer that the strict branch succeeds; it can only be used to analyze the open-system bypass companion.

\begin{theorem}[Part VIII relaxed-branch and export-companion theorem]
\label{thm:part08-main}
Inside the declared relaxed-constraint branch, the one-mode leakage/work closure, the non-rigid mouth/dressing closure, the compensated two-mode source closure, the reduced one-dimensional event-chain closure, the characteristic crossing/collapse closure, the microscopic cubic/quintic export closure, and the physical-calibration companion, the following statements hold.
\begin{enumerate}[leftmargin=2.2em]
  \item The relaxed branch is a codimension-three lift of the Stage 242 front end with exact recovery map
  \begin{equation}
  \label{eq:part08-main-recovery-map}
  \ell_w=0,
  \qquad
  f_U=0,
  \qquad
  a=b=0.
  \end{equation}
  On this slice all lifted leakage, work, non-rigid, and compensated-source observables reduce to the strict standard-recovery values.
  \item The same-charge static and linear dynamic kernel span remains short-ranged:
  \begin{equation}
  \label{eq:part08-main-short-range-span}
  \mathcal K_{\rm SR}
  =
  \operatorname{span}\left\{
  r^{-6},
  \frac{e^{-2\kappa r}}{r^4},
  \frac{e^{-4\kappa r}}{r^2}
  \right\}.
  \end{equation}
  In particular,
  \begin{equation}
  \label{eq:part08-main-short-range-limits}
  \lim_{r\to\infty}r\,\delta V_{\rm stat}(r)=0,
  \qquad
  \lim_{r\to\infty}r\,\Re\mathfrak V_{\rm dyn}(r,\omega)=0.
  \end{equation}
  \item On the selected-support branch, the leakage/work lane is explicit and support-side:
  \begin{equation}
  \label{eq:part08-main-leak-support}
  S_{\rm leak}(r)
  =
  \frac{8\sqrt2\,\eta_{\rm leak}\mu_w\rho_0}{\pi^{5/2}\lambda^3}
  \Lambda(r)\varrho(r),
  \end{equation}
  \begin{equation}
  \label{eq:part08-main-work-support}
  \mathcal W_w^{\rm sess}(r)
  =
  \frac{512\,\eta_{\rm leak}^{2}\mu_w q\rho_0}{\pi^4\lambda^2}
  \Lambda(r)^2\varrho(r)^2.
  \end{equation}
  These quantities are independent of the orbit-lock variables \((R_{\rm tr},R_{\rm target},\epsilon_\eta)\) within the declared pullback.
  \item The non-rigid mouth packet is exact in the physical logarithmic chart
  \begin{equation}
  \label{eq:part08-main-uv-chart}
  U=\ln\frac{\mathcal T^2}{\mathcal T_{\rm ref}^2},
  \qquad
  V=\ln\frac{\epsilon_\eta}{\epsilon_{\eta,{\rm ref}}}.
  \end{equation}
  With
  \begin{equation}
  \label{eq:part08-main-delta-uv}
  \Delta_{UV}=a_Ua_V-\chi_{UV}^{2},
  \end{equation}
  the stationary point is
  \begin{equation}
  \label{eq:part08-main-uv-stationary}
  U=\frac{a_Vf_U}{\Delta_{UV}},
  \qquad
  V=\frac{\chi_{UV}f_U}{\Delta_{UV}},
  \qquad
  \Delta_{UV}>0.
  \end{equation}
  The corresponding drain is nonnegative:
  \begin{equation}
  \label{eq:part08-main-uv-drain}
  \mathcal D_{UV}=\frac{\chi_{UV}^{2}a_Vf_U^2}{\Delta_{UV}^{2}}\ge 0.
  \end{equation}
  \item The first compensated two-mode source family is an exact two-moment compiler:
  \begin{equation}
  \label{eq:part08-main-source-family}
  \sigma_{a,b}(x)=1+a\cos(\pi x)+b\cos(2\pi x),
  \qquad
  \int_0^1\sigma_{a,b}(x)\,dx=1,
  \end{equation}
  with
  \begin{equation}
  \label{eq:part08-main-source-moments}
  \mathfrak g(a,b)=\frac{2}{\pi}\left(1+\frac{a}{3}-\frac{b}{15}\right),
  \qquad
  \mathcal S(a,b)=\frac{2\tanh(\pi/2)}{\pi}\left(1+\frac{a}{5}+\frac{b}{17}\right).
  \end{equation}
  Its normalized two-moment map has determinant \(14/425>0\), so the first compensated source family spans the two carried source moments exactly.
  \item The relaxed stationary barrier is exactly \cref{eq:part08-intro-veff230}, and the lowering identity is
  \begin{equation}
  \label{eq:part08-main-lowering-identity}
  V_{\rm short}^{(1p)}(r)-V_{\rm eff}^{(247)}(r)
  =
  \lambda_LS_{\rm leak}(r)
  +\lambda_W\mathcal W_w^{\rm sess}(r)
  +\Delta E_{UV}(r)
  +\mathcal M_\sigma(r).
  \end{equation}
  Hence, when the imported packet scalars and weights are nonnegative, \(V_{\rm eff}^{(247)}\le V_{\rm short}^{(1p)}\) pointwise.
  \item The dynamic event chain is determined by the conserved energy \cref{eq:part08-intro-event-energy}.  The lowered-branch classical threshold speed is
  \begin{equation}
  \label{eq:part08-main-vcrit-new}
  v_{\rm crit,new}=\sqrt{\frac{2}{m_s}\bigl(V_{\rm peak}-V(r_0)\bigr)},
  \end{equation}
  while the pure-Coulomb contact threshold is
  \begin{equation}
  \label{eq:part08-main-vcrit-coul}
  v_{\rm contact,Coul}
  =
  \sqrt{\frac{2}{m_s}\left(\frac1{r_{\rm contact}}-\frac1{r_0}\right)}.
  \end{equation}
  If \(v_{\rm crit,new}<v_0<v_{\rm contact,Coul}\), the lowered branch reaches contact while the pure Coulomb branch still turns back.
  \item The subbarrier WKB action is
  \begin{equation}
  \label{eq:part08-main-wkb-action}
  I_{\rm new}(E)=
  \frac{1}{\hbar_{\rm eff}}
  \int_{r_-(E)}^{r_+(E)}
  \sqrt{2m_s\bigl(V(r)-E\bigr)}\,dr,
  \end{equation}
  and the transmission ratio relative to Coulomb is
  \begin{equation}
  \label{eq:part08-main-transmission-ratio}
  \frac{T_{\rm new}(E)}{T_{\rm Coul}(E)}
  =
  \exp\bigl[-2\bigl(I_{\rm new}(E)-I_{\rm Coul}(E)\bigr)\bigr].
  \end{equation}
  \item The helicity-export comparison is diagnostic.  With orientation label \(\sigma=\pm1\), the reduced closure gives
  \begin{equation}
  \label{eq:part08-main-helicity-rate}
  \dot H_{{\rm sub},\sigma}=\Gamma_0(1+\sigma\alpha_h),
  \qquad
  |\alpha_h|<1
  \end{equation}
  for positive export in both branches.  The aligned branch dominates whenever \(\alpha_h>0\), but this does not by itself prove a microscopic spin sector.
  \item The crossing-versus-collapse safe edge is
  \begin{equation}
  \label{eq:part08-main-goldilocks-edge}
  E_{\rm edge}=V_{\rm peak}
  +
  \frac{m_s}{\mu_\eta}
  \frac{g_{UV}\chi_{\rm peak}\lambda_{\rm eff}^{2}}{2},
  \end{equation}
  equivalently
  \begin{equation}
  \label{eq:part08-main-goldilocks-speed}
  v_{\rm safe,min}^{2}
  =
  v_{\rm crit,new}^{2}
  +
  \frac{\lambda_{\rm eff}^{2}g_{UV}\chi_{\rm peak}}{\mu_\eta}.
  \end{equation}
  Under the declared characteristic closure this is a one-sided lower edge.
  \item The microscopic active-leg export kernel is super-Ohmic:
  \begin{equation}
  \label{eq:part08-main-kexp}
  K_{\rm exp}(s)=\Gamma_3s^3+\Gamma_5s^5,
  \qquad
  \Gamma_3,\Gamma_5\ge0.
  \end{equation}
  The positive growth root is reduced but not eliminated by any finite passive \((\Gamma_3,\Gamma_5)\) in this minimal closure.  The exact event-safe condition is
  \begin{equation}
  \label{eq:part08-main-safe-half-plane}
  \widehat\Gamma_3+s_c^2\widehat\Gamma_5
  \ge
  \frac{s_0^2-s_c^2}{s_c^3},
  \qquad
  \widehat\Gamma_n:=\Gamma_n/\mu_\eta.
  \end{equation}
  \item The heat partition is channel-resolved.  For event-shape quotient \(\mathfrak r_V=\mathcal I_3/\mathcal I_2\),
  \begin{equation}
  \label{eq:part08-main-heat-fraction}
  f_{\rm lat}(\mathfrak r_V)=
  \frac{\Gamma_3^{\rm lat}+\Gamma_5^{\rm lat}\mathfrak r_V}{\Gamma_3+\Gamma_5\mathfrak r_V},
  \end{equation}
  and on a single-growth cold event \(\mathfrak r_V=s^2\).  At the safe edge, the total minimum exported energy is
  \begin{equation}
  \label{eq:part08-main-safe-energy}
  E_{\rm exp,min}^{\rm safe}
  =
  \frac{V_{\rm in}^{2}}{2}(e^2-1)\mu_\eta(s_0^2-s_c^2).
  \end{equation}
  \item The physical-material companion is an exact calibration image of the Stage 252 event-equivalent rate.  With calibration factor \(\Upsilon_{\rm lat}>0\),
  \begin{equation}
  \label{eq:part08-main-ep-threshold}
  (\lambda_{\rm ep}\omega_D)_{\min}^{(253)}
  =
  \frac{f_{\rm lat}(s_c)\mu_\eta(s_0^2-s_c^2)}
  {\Upsilon_{\rm lat}\zeta_{\rm ep}s_ct_*}.
  \end{equation}
  Candidate hosts are screened by \cref{eq:part08-intro-material-screen}.
\end{enumerate}
Every item involving actual branch data, microscopic coefficients, physical units, or candidate material constants remains a realization or calibration question beyond the algebraic compiler.
\end{theorem}

\section{Relaxed constraint branch declaration}
\label{sec:part08_relaxed_branch:relaxed-constraint-branch-declaration}

Stage 243 declares the relaxed branch without yet assembling a barrier.  Its purpose is to make the Session-I relaxations legal inside the derivation tree while preserving the same-charge short-range verdict.  The inherited base object is the Stage 242 front-end packet
\begin{equation}
\label{eq:part08-p225-packet}
\mathcal P_{242}
=
\Bigl(
\epsilon,
\varrho_{\rm phys},
\sigma_{\rm phys},
\text{ranking region},
R_{\rm tr},
R_{\rm target},
\epsilon_\eta,
 d\ln R_{\rm tr},
 d\ln R_{\rm target},
 d\ln\epsilon_\eta,
 N_Q-1
\Bigr),
\end{equation}
with selected-support packet
\begin{equation}
\label{eq:part08-s225-packet}
\mathcal S_{242}=(\zeta,M_{\rm mix},S(\zeta;\epsilon),M_{\rm tr}).
\end{equation}
The strict slice is
\begin{equation}
\label{eq:part08-std-branch-slice}
\mathfrak B_{\rm std}
=
\left\{
J^w=0,
\;U=0,
\;V=0,
\;\varsigma(z)\equiv 1
\right\}
\subset(\mathcal P_{242},\mathcal S_{242}).
\end{equation}
This definition is useful because it names exactly what is being relaxed: transverse-current suppression, rigid-mouth orbit lock, and positive-source Family-1 closure.

\subsection{Open-system leakage/work lane}
\label{subsec:part08-open-lane}

The open-system lift begins with the exact projected-continuity leakage identity
\begin{equation}
\label{eq:part08-stage243-leakage-identity}
S_{\rm leak}
=
-\bigl[W(w)j^w\bigr]_{-\infty}^{+\infty}
+
\int_{-\infty}^{+\infty}W'(w)j^w(w)\,dw.
\end{equation}
For the declaration-level Gaussian closure,
\begin{equation}
\label{eq:part08-stage243-gaussian-profile}
W(w)=\frac{e^{-w^2}}{\sqrt\pi},
\qquad
j^w(w)=\ell_w j_0\,w e^{-w^2},
\qquad
E_w(w)=E_0we^{-w^2},
\end{equation}
so the boundary term vanishes and
\begin{equation}
\label{eq:part08-stage243-leak-work}
S_{\rm leak}=-\frac{\sqrt2}{4}\ell_wj_0,
\qquad
\mathcal W_w=\frac{\sqrt{2\pi}}{8}\ell_wj_0E_0.
\end{equation}
The declaration packet for this lane is
\begin{equation}
\label{eq:part08-stage243-lw}
L_w=(S_{\rm leak},\mathcal W_w).
\end{equation}
Both components vanish on \(\ell_w=0\).

\subsection{Non-rigid mouth/dressing lane}
\label{subsec:part08-uv-lane}

The non-rigid lift uses the minimal free energy
\begin{equation}
\label{eq:part08-stage243-fuv}
\mathcal F_{UV}(U,V)
=
\frac12k_UU^2+\frac12k_VV^2-\chi_\lambda UV-f_UU.
\end{equation}
Stationarity gives
\begin{equation}
\label{eq:part08-stage243-uv-solution}
U=\frac{k_Vf_U}{k_Uk_V-\chi_\lambda^2},
\qquad
V=\frac{\chi_\lambda f_U}{k_Uk_V-\chi_\lambda^2},
\qquad
\frac{V}{U}=\frac{\chi_\lambda}{k_V}.
\end{equation}
The branch is admissible when
\begin{equation}
\label{eq:part08-stage243-uv-stability}
k_Uk_V-\chi_\lambda^2>0,
\end{equation}
and the drain is
\begin{equation}
\label{eq:part08-stage243-duv}
\mathcal D_{UV}=\chi_\lambda UV
=
\frac{\chi_\lambda^2k_Vf_U^2}{(k_Uk_V-\chi_\lambda^2)^2}\ge0.
\end{equation}
This lane is the finite relaxed continuation of the earlier weak-axisymmetric transfer-shape scalar: at infinitesimal order, \(\Xi_1=\delta U\).

\subsection{Compensated sign-changing source lane}
\label{subsec:part08-compensated-source-lane}

The compensated source lift is
\begin{equation}
\label{eq:part08-stage243-varsigma}
\varsigma(z)=1+a\cos(\pi z)+b\cos(2\pi z),
\qquad
z\in[0,1],
\end{equation}
with exact unit mean.  Setting \(y=\cos(\pi z)\),
\begin{equation}
\label{eq:part08-stage243-varsigma-quadratic}
\varsigma(y)=1-b+ay+2by^2,
\qquad
-1\le y\le 1.
\end{equation}
The interior stationary point and value are
\begin{equation}
\label{eq:part08-stage243-varsigma-min}
y_*=-\frac{a}{4b},
\qquad
\varsigma(y_*)=1-b-\frac{a^2}{8b},
\end{equation}
with the boundary candidates \(1+a+b\) and \(1-a+b\).  Thus the first source relaxation is not a vague sign-changing profile; it is an exact mean-preserving quadratic problem.

The full declaration packet is
\begin{equation}
\label{eq:part08-stage243-relaxed-branch}
\mathfrak B_{243}^{\rm relax}
=
\left\{(\mathcal P_{242},\mathcal S_{242});\;L_w,\;L_{UV},\;L_\varsigma\right\},
\end{equation}
where
\begin{equation}
\label{eq:part08-stage243-lanes}
L_w=(S_{\rm leak},\mathcal W_w),
\quad
L_{UV}=(U,V,\mathcal D_{UV}),
\quad
L_\varsigma=(\varsigma,\varsigma_{\min},\mathrm{signchg}).
\end{equation}

\subsection{Short-range theorem for the lifted branch}
\label{subsec:part08-short-range-theorem}

The most important Stage 243 output is the short-range firewall.  With primitive profiles
\begin{equation}
\label{eq:part08-stage243-source-profiles}
\mathcal S_Q(r)=r^{-3},
\qquad
\mathcal S_Y(r)=\frac{e^{-2\kappa r}}{r},
\end{equation}
the static one-port correction is
\begin{equation}
\label{eq:part08-stage243-static-correction}
\delta V_{\rm stat}(r)
=
-\frac12\left[
\frac{\mathcal C_6}{r^6}
+2\mathcal C_4\frac{e^{-2\kappa r}}{r^4}
+\mathcal C_2\frac{e^{-4\kappa r}}{r^2}
\right],
\end{equation}
and the linear monochromatic dynamic bundle has the same spatial span:
\begin{equation}
\label{eq:part08-stage243-dynamic-correction}
\Re\mathfrak V_{\rm dyn}(r,\omega)
=
-\frac12\left[
\frac{\mathcal C_6(\omega)}{r^6}
+2\mathcal C_4(\omega)\frac{e^{-2\kappa r}}{r^4}
+\mathcal C_2(\omega)\frac{e^{-4\kappa r}}{r^2}
\right].
\end{equation}
Thus the relaxed branch is a short-range/open-system bypass branch.  It can reshape near-contact energetics, but it cannot be cited as a new long-range same-charge force.

\section{Selected leakage and scalar-photon work compiler}
\label{sec:part08_relaxed_branch:selected-leakage-and-scalar-photon-work-compiler}

Stage 244 takes the declared \(J^w\neq0\) lane and attaches it to the selected-support packet.  The parent theory supplies both exact identities needed for the compiler: projected continuity gives \cref{eq:part08-stage243-leakage-identity}, and the localized Maxwell energy ledger gives
\begin{equation}
\label{eq:part08-stage244-em-energy}
\partial_tu_{\rm EM}+\partial_AS_{\rm EM}^{A}
=
-J^AE_A
=-(J^aE_a+J^wE_w).
\end{equation}
Thus \(S_{\rm leak}\) and \(J^wE_w\) are legitimate reduced observables once the mixed lane is opened.

\subsection{Gaussian one-mode leakage law}
\label{subsec:part08-stage244-gaussian-leakage}

The Stage 244 one-mode closure uses
\begin{equation}
\label{eq:part08-stage244-wlambda}
W_\lambda(w)=\frac{e^{-w^2/\lambda^2}}{\lambda\sqrt\pi},
\qquad
\phi_\lambda(w)=\frac{2w}{\sqrt\pi\lambda^3}e^{-w^2/\lambda^2},
\end{equation}
with
\begin{equation}
\label{eq:part08-stage244-ew-jw}
E_w(w;r)=-\mathcal E_0(r)\phi_\lambda(w),
\qquad
j^w(w;r)=\mu_w\rho_0E_w(w;r),
\qquad
J^w=qj^w.
\end{equation}
The exact projected leakage is
\begin{equation}
\label{eq:part08-stage244-sleak-e0}
S_{\rm leak}(r)
=
\frac{\sqrt2\mu_w\rho_0}{2\sqrt\pi\lambda^3}\mathcal E_0(r).
\end{equation}
The exact one-mode bulk work scalar is
\begin{equation}
\label{eq:part08-stage244-bulk-work}
\mathcal W_w^{\rm bulk}(r)
=
q\int j^wE_w\,dw
=
\frac{\sqrt2\mu_wq\rho_0}{2\sqrt\pi\lambda^3}\mathcal E_0(r)^2,
\end{equation}
so
\begin{equation}
\label{eq:part08-stage244-work-leak-relation}
\mathcal W_w^{\rm bulk}(r)=q\mathcal E_0(r)S_{\rm leak}(r).
\end{equation}
The Session-I scalar-photon work variable is the rescaled scalar
\begin{equation}
\label{eq:part08-stage244-session-work}
\mathcal W_w^{\rm sess}(r)
=
\frac{2\mu_wq\rho_0}{\lambda^2}\mathcal E_0(r)^2
=
2\sqrt{2\pi}\lambda\,\mathcal W_w^{\rm bulk}(r),
\end{equation}
or equivalently
\begin{equation}
\label{eq:part08-stage244-quadratic-law}
\mathcal W_w^{\rm sess}(r)
=
\frac{4\pi q\lambda^4}{\mu_w\rho_0}S_{\rm leak}(r)^2.
\end{equation}

\subsection{Selected-support pullback}
\label{subsec:part08-stage244-selected-pullback}

The selected-support demand from Stage 242 is
\begin{equation}
\label{eq:part08-stage244-pitr}
\Pi_{\rm tr}
=
\frac43C_{\rm mix}
=
\frac{32\Lambda(1-
\epsilon)}{3\pi^2}
=
\frac{16}{\pi^2}\Lambda\varrho,
\qquad
\varrho=\frac23(1-\epsilon).
\end{equation}
Stage 244 chooses the simplest amplitude pullback
\begin{equation}
\label{eq:part08-stage244-e0-pullback}
\mathcal E_0(r)=\eta_{\rm leak}\Pi_{\rm tr}(r).
\end{equation}
Substitution gives the selected-branch leakage law \cref{eq:part08-main-leak-support} and the selected-branch work law \cref{eq:part08-main-work-support}.  In the \((\Lambda,\epsilon)\) notation the same work law is
\begin{equation}
\label{eq:part08-stage244-work-epsilon}
\mathcal W_w^{\rm sess}(r)
=
\frac{2048\eta_{\rm leak}^{2}\mu_wq\rho_0}{9\pi^4\lambda^2}
\Lambda(r)^2(1-\epsilon(r))^2.
\end{equation}
The formulas depend only on selected-support variables.  Therefore,
\begin{equation}
\label{eq:part08-stage244-support-orbit-split}
\partial_{R_{\rm tr}}S_{\rm leak}
=
\partial_{R_{\rm target}}S_{\rm leak}
=
\partial_{\epsilon_\eta}S_{\rm leak}=0,
\end{equation}
and the same statement holds for both work scalars within the declared pullback.

The orientation parity is exact:
\begin{equation}
\label{eq:part08-stage244-parity}
S_{\rm leak}(-\eta_{\rm leak})=-S_{\rm leak}(\eta_{\rm leak}),
\qquad
\mathcal W_w(-\eta_{\rm leak})=\mathcal W_w(\eta_{\rm leak}).
\end{equation}
Thus the leakage sign is orientation-sensitive, but the exported-work magnitude is not.

\section{Non-rigid mouth and \texorpdfstring{\((U,V)\)}{(U,V)} drain packet}
\label{sec:part08_relaxed_branch:non-rigid-mouth-and-u-v-drain-packet}

Stage 245 compiles the abstract non-rigid lane into the actual finite physical variables used by the rigid-mouth chart.  The carried chart is \cref{eq:part08-main-uv-chart}; hence
\begin{equation}
\label{eq:part08-stage245-direct-observables}
\mathcal T^2=\mathcal T_{\rm ref}^{2}e^U,
\qquad
\epsilon_\eta=\epsilon_{\eta,{\rm ref}}e^V.
\end{equation}
The selected-branch identity
\begin{equation}
\label{eq:part08-stage245-selected-identity}
R_{\rm target}\mathcal T^2=\Lambda_0(1-\epsilon_\eta)
\end{equation}
then gives
\begin{equation}
\label{eq:part08-stage245-rtarget-finite}
\frac{R_{\rm target}}{R_{\rm target,ref}}
=
\frac{1-\epsilon_{\eta,{\rm ref}}e^V}{1-\epsilon_{\eta,{\rm ref}}}e^{-U}.
\end{equation}
So the finite non-rigid packet has two independent legs: transfer shape and dressing.

The reduced free energy in session notation is
\begin{equation}
\label{eq:part08-stage245-fnr}
\mathcal F_{\rm nr}(U,V;r)
=
\frac12a_UU^2+\frac12a_VV^2-\chi_{UV}(r)UV-f_U(r)U.
\end{equation}
The stationary solution is \cref{eq:part08-main-uv-stationary}, with \(\Delta_{UV}\) given by \cref{eq:part08-main-delta-uv}.  In finite physical variables,
\begin{equation}
\label{eq:part08-stage245-finite-t}
\mathcal T^2(r)
=
\mathcal T_{\rm ref}^2\exp\left(\frac{a_Vf_U(r)}{\Delta_{UV}(r)}\right),
\end{equation}
\begin{equation}
\label{eq:part08-stage245-finite-epseta}
\epsilon_\eta(r)
=
\epsilon_{\eta,{\rm ref}}\exp\left(\frac{\chi_{UV}(r)f_U(r)}{\Delta_{UV}(r)}\right),
\end{equation}
\begin{equation}
\label{eq:part08-stage245-finite-rtarget}
\frac{R_{\rm target}(r)}{R_{\rm target,ref}}
=
\frac{1-\epsilon_{\eta,{\rm ref}}\exp(\chi_{UV}f_U/\Delta_{UV})}{1-\epsilon_{\eta,{\rm ref}}}
\exp(-a_Vf_U/\Delta_{UV}).
\end{equation}

The dependent microscopic correction inherited from the rigid-mouth plane is
\begin{equation}
\label{eq:part08-stage245-dependent-plane}
(\Delta_T,\Delta_{K_\eta},\Delta_\mu)
=(0,-V,U-V),
\end{equation}
so the non-rigid stationary packet induces
\begin{equation}
\label{eq:part08-stage245-dependent-nr}
\mathbf y_{\rm nr}^{\rm dep}(r)
=
\begin{pmatrix}
0\\[2pt]
-\chi_{UV}f_U/\Delta_{UV}\\[2pt]
(a_V-\chi_{UV})f_U/\Delta_{UV}
\end{pmatrix}.
\end{equation}
The drain \(\mathcal D_{UV}\) is given by \cref{eq:part08-main-uv-drain}; it is nonnegative even when \(U\) and \(V\) have opposite signs.

At weak-axisymmetric order, write
\begin{equation}
\label{eq:part08-stage245-first-order-uv}
U=\varepsilon u_1+O(\varepsilon^2),
\qquad
V=\varepsilon v_1+O(\varepsilon^2).
\end{equation}
Then
\begin{equation}
\label{eq:part08-stage245-xi-r-first}
\Xi_1^{\rm nr}=u_1,
\qquad
\mathcal R_1^{\rm nr}+\Xi_1^{\rm nr}
=-\frac{\epsilon_{\eta,{\rm ref}}}{1-\epsilon_{\eta,{\rm ref}}}v_1,
\end{equation}
\begin{equation}
\label{eq:part08-stage245-r-first}
\mathcal R_1^{\rm nr}
=
-u_1-
\frac{\epsilon_{\eta,{\rm ref}}}{1-\epsilon_{\eta,{\rm ref}}}v_1.
\end{equation}
Thus the direct outgoing-prefactor scalar remains the \(U\) leg, while active dressing adds an independent selected-branch residual.

Within the declared closure, this packet is orbit-side/support-blind:
\begin{equation}
\label{eq:part08-stage245-support-blind}
\partial_\Lambda U=
\partial_\varrho U=
\partial_\Lambda V=
\partial_\varrho V=0,
\end{equation}
provided \(f_U\) and \(\chi_{UV}\) are treated as orbit-side data rather than support-placement coordinates.  This complements the support-side result from Stage 244.

\section{Compensated multimode source compiler}
\label{sec:part08_relaxed_branch:compensated-multimode-source-compiler}

Stage 246 compiles the source lane that Stage 243 only declared.  The carried source kernels on the dimensionless mouth coordinate \(x\in[0,1]\) are
\begin{equation}
\label{eq:part08-stage246-source-kernels}
c(x)=\cos\frac{\pi x}{2},
\qquad
K_q(x)=\frac{\cosh\left(\frac\pi2(1-x)\right)}{\cosh(\pi/2)}.
\end{equation}
The two source moments are
\begin{equation}
\label{eq:part08-stage246-source-moments-def}
\mathfrak g[\sigma]=\int_0^1\sigma(x)c(x)\,dx,
\qquad
\mathcal S[\sigma]=\int_0^1\sigma(x)K_q(x)\,dx.
\end{equation}
The carried Family-1 mixed-to-shell loading ratio is
\begin{equation}
\label{eq:part08-stage246-rsigma-def}
\mathcal R[\sigma]
=
\frac{(\mathfrak g[\sigma]-\mathfrak r_{F1})^2}{1+\mathfrak r_{F1}^2},
\qquad
\mathfrak r_{F1}=
\sqrt{\frac{12}{\pi^2}\left(\frac{37}{20}\right)^2-1}.
\end{equation}
Stage 246 keeps this law and extends the source family beyond the positive corridor.

The compensated family is \cref{eq:part08-main-source-family}.  Its exact minimum is
\begin{equation}
\label{eq:part08-stage246-sigma-min}
\sigma_{\min}(a,b)=
\begin{cases}
1+b-|a|, & b\le0\;\text{or}\; |a|\ge4b,\\[4pt]
1-b-\dfrac{a^2}{8b}, & b>0\;\text{and}\; |a|\le4b.
\end{cases}
\end{equation}
Thus
\begin{equation}
\label{eq:part08-stage246-sign-change}
\mathrm{signchg}(a,b)\iff \sigma_{\min}(a,b)<0.
\end{equation}
The exact moment formulas are \cref{eq:part08-main-source-moments}.  In normalized variables
\begin{equation}
\label{eq:part08-stage246-normalized-source-moments}
\widetilde g:=\frac\pi2\mathfrak g-1,
\qquad
\widetilde S:=\frac{\pi}{2\tanh(\pi/2)}\mathcal S-1,
\end{equation}
the two-moment map is
\begin{equation}
\label{eq:part08-stage246-matrix}
\begin{pmatrix}
\widetilde g\\[2pt]
\widetilde S
\end{pmatrix}
=
\begin{pmatrix}
1/3 & -1/15\\[4pt]
1/5 & 1/17
\end{pmatrix}
\begin{pmatrix}
a\\[2pt]
b
\end{pmatrix},
\qquad
\det M_{\rm src}=\frac{14}{425}.
\end{equation}
The inverse is
\begin{equation}
\label{eq:part08-stage246-inverse-source-map}
a=\frac{85}{42}\widetilde S+\frac{25}{14}\widetilde g,
\qquad
b=\frac{425}{42}\widetilde S-\frac{85}{14}\widetilde g.
\end{equation}

For the transported branch,
\begin{equation}
\label{eq:part08-stage246-transport}
a(r)=a_0s(r),
\qquad
b(r)=b_0s(r),
\qquad
s(r)=\frac{r_\sigma^2}{r^2+r_\sigma^2}.
\end{equation}
The output packet is
\begin{equation}
\label{eq:part08-stage246-msigma-packet}
\mathcal M_\sigma(r)=
\bigl(
\sigma(x;r),
\sigma_{\min}(r),
\mathrm{signchg}(r),
\mathfrak g[\sigma](r),
\mathcal S[\sigma](r),
\mathcal R[\sigma](r)
\bigr),
\end{equation}
with
\begin{equation}
\label{eq:part08-stage246-transported-moments}
\mathfrak g[\sigma](r)
=
\frac{2}{\pi}\left(1+\frac{a(r)}{3}-\frac{b(r)}{15}\right),
\qquad
\mathcal S[\sigma](r)
=
\frac{2\tanh(\pi/2)}{\pi}\left(1+\frac{a(r)}{5}+\frac{b(r)}{17}\right).
\end{equation}
On the Session-I orientation \(a_0>0\), \(b_0<0\), the minimum is boundary-controlled:
\begin{equation}
\label{eq:part08-stage246-session-sign-threshold}
\sigma_{\min}(r)=1-(a_0-b_0)s(r),
\qquad
\mathrm{signchg}(r)\iff r<r_\sigma\sqrt{a_0-b_0-1}.
\end{equation}
This is the source-side packet inserted into the stationary barrier compiler.

\section{Relaxed stationary barrier front end}
\label{sec:part08_relaxed_branch:relaxed-stationary-barrier-front-end}

Stage 247 assembles the first stationary lowered barrier from the one-port short-range baseline and the three relaxed packets.  The one-port static bundle is controlled by
\begin{equation}
\label{eq:part08-stage247-delta-q-p}
\Delta=\Omega_U^2\Omega_W^2-R_{\rm mix}^{2},
\qquad
Q=G_U^2\Omega_W^2+2G_UG_WR_{\rm mix}+G_W^2\Omega_U^2,
\end{equation}
\begin{equation}
\label{eq:part08-stage247-p-d0}
P=\Omega_U^2G_W+R_{\rm mix}G_U,
\qquad
D_0=K_*-\frac{Q}{\Delta}.
\end{equation}
With primitive short-range source amplitudes \(\beta_Q,\beta_U,\beta_W\), the product-family coefficients are
\begin{equation}
\label{eq:part08-stage247-c6c4c2}
\mathcal C_6=\chi_{qq}\beta_Q^2,
\qquad
\mathcal C_4=\chi_{qU}\beta_Q\beta_U+\chi_{qW}\beta_Q\beta_W,
\end{equation}
\begin{equation}
\label{eq:part08-stage247-c2}
\mathcal C_2=\chi_{UU}\beta_U^2+2\chi_{UW}\beta_U\beta_W+\chi_{WW}\beta_W^2.
\end{equation}
The collected baseline is
\begin{equation}
\label{eq:part08-stage247-vshort}
V_{\rm short}^{(1p)}(r)
=
\frac1r\left(1+\frac12e^{-2\kappa r}\right)
-
\frac{A_6}{r^6}
-
A_4\frac{e^{-2\kappa r}}{r^4}
-
A_2\frac{e^{-4\kappa r}}{r^2},
\end{equation}
where
\begin{equation}
\label{eq:part08-stage247-a6a4a2}
A_6=3\alpha_6+\frac12\mathcal C_6,
\qquad
A_4=\mathcal C_4,
\qquad
A_2=\alpha_2+\frac12\mathcal C_2.
\end{equation}
This baseline inherits the same short-range firewall from Stages 219--243.

The imported relaxed scalars are
\begin{equation}
\label{eq:part08-stage247-lag-l}
\mathfrak L(r):=\Lambda(r)\varrho(r),
\qquad
S_{\rm leak}=\frac{8\sqrt2\eta_{\rm leak}\mu_w\rho_0}{\pi^{5/2}\lambda^3}\mathfrak L,
\end{equation}
\begin{equation}
\label{eq:part08-stage247-work}
\mathcal W_w^{\rm sess}
=
\frac{512\eta_{\rm leak}^{2}\mu_wq\rho_0}{\pi^4\lambda^2}\mathfrak L^2,
\qquad
\Delta E_{UV}=\eta_{UV}\mathcal D_{UV},
\end{equation}
\begin{equation}
\label{eq:part08-stage247-source-lowering}
\mathcal M_\sigma(r)=\xi_R\bigl[\mathcal R_\infty-\mathcal R(r)\bigr],
\qquad
\mathcal R_\infty=\frac{(2/\pi-\mathfrak r_{F1})^2}{1+\mathfrak r_{F1}^2}.
\end{equation}
The stationary compiler is \cref{eq:part08-intro-veff230}; the lowering identity is \cref{eq:part08-main-lowering-identity}.

The Session-I benchmark is useful as an audit point, not as a general theorem.  At the strongest softening point, the Stage 247 decomposition reads
\begin{equation}
\label{eq:part08-stage247-benchmark-decomposition}
3.74163698
-
0.26971918\times0.31069599
-
1.51632107
-
0.21064278
-
0.18386120
=
1.74701126.
\end{equation}
This shows that, once the one-port baseline is fixed, the recorded stationary lowering can be decomposed into work, \(U/V\) drain, compensated source response, and one remaining leakage-to-barrier embedding coefficient.  It does not change the short-range kernel class.

\section{Dynamic event-chain compiler}
\label{sec:part08_relaxed_branch:dynamic-event-chain-compiler}

Stage 248 promotes \(V_{\rm eff}^{(247)}\) into a reduced scattering/event-chain compiler.  Write
\begin{equation}
\label{eq:part08-stage248-vdef}
V(r):=V_{\rm eff}^{(247)}(r).
\end{equation}
The reduced radial equation is
\begin{equation}
\label{eq:part08-stage248-radial-eom}
m_s\ddot r=-V'(r),
\end{equation}
with conserved energy \cref{eq:part08-intro-event-energy}.  If the branch has a single barrier peak
\begin{equation}
\label{eq:part08-stage248-peak}
V'(r_{\rm peak})=0,
\qquad
V''(r_{\rm peak})<0,
\qquad
V_{\rm peak}=V(r_{\rm peak}),
\end{equation}
then the finite-radius threshold law is \cref{eq:part08-main-vcrit-new}.  The pure-Coulomb contact threshold is \cref{eq:part08-main-vcrit-coul}.  The classical comparison window is therefore
\begin{equation}
\label{eq:part08-stage248-classical-window}
v_{\rm crit,new}<v_0<v_{\rm contact,Coul}.
\end{equation}

For subbarrier energy \(V(r_0)<E<V_{\rm peak}\), turning points solve
\begin{equation}
\label{eq:part08-stage248-turning-points}
V(r_-(E))=E,
\qquad
V(r_+(E))=E,
\qquad
r_-(E)<r_+(E),
\end{equation}
and transport as
\begin{equation}
\label{eq:part08-stage248-turning-transport}
\frac{dr_\pm}{dE}=\frac1{V'(r_\pm(E))}.
\end{equation}
The WKB action and transmission ratio are \cref{eq:part08-main-wkb-action,eq:part08-main-transmission-ratio}.  The pure-Coulomb outer turning point is exact,
\begin{equation}
\label{eq:part08-stage248-coul-turn}
r_{\rm turn,Coul}(E)=\frac1E,
\end{equation}
and the Coulomb action closes as
\begin{equation}
\label{eq:part08-stage248-coul-action}
I_{\rm Coul}(E)
=
\frac{\sqrt{2m_s}}{\hbar_{\rm eff}}
\left[
\frac{\pi}{2\sqrt E}
-
\sqrt{r_{\rm contact}(1-Er_{\rm contact})}
-
\frac{\arcsin\sqrt{Er_{\rm contact}}}{\sqrt E}
\right]
\end{equation}
for \(Er_{\rm contact}<1\).

Near the top, with
\begin{equation}
\label{eq:part08-stage248-near-top}
V(r)=V_{\rm peak}-\frac{K_{\rm peak}}2(r-r_{\rm peak})^2+O((r-r_{\rm peak})^3),
\qquad
K_{\rm peak}=-V''(r_{\rm peak}),
\end{equation}
the leading top action is
\begin{equation}
\label{eq:part08-stage248-top-action}
I_{\rm top}(E)=
\frac{\pi(V_{\rm peak}-E)}{\hbar_{\rm eff}}
\sqrt{\frac{m_s}{K_{\rm peak}}}
+O((V_{\rm peak}-E)^{3/2}).
\end{equation}
The dynamic diagnostics carried forward are
\begin{equation}
\label{eq:part08-stage248-diagnostics}
\Xi_{\rm turn}(E)=\Xi_1(r_+(E)),
\qquad
\lambda_{\rm th}(E)=\left|\frac{E}{V'(r_+(E))}\right|.
\end{equation}
These tags are later used in the helicity and Goldilocks compilers; they do not alter the scalar event chain.

\section{Helicity export diagnostic}
\label{sec:part08_relaxed_branch:helicity-export-diagnostic}

Stage 249 attaches a hidden-channel export diagnostic to the Stage 248 event chain.  It is diagnostic rather than constitutive.  The parent projected helicity ledger gives the exact subscale helicity
\begin{equation}
\label{eq:part08-stage249-hsub}
h_{\rm sub}
=
\overline h-h_{\rm res}
=
\overline{\mathbf A'\cdot\mathbf B'},
\end{equation}
with transfer equation
\begin{equation}
\label{eq:part08-stage249-hsub-eq}
\partial_th_{\rm sub}+\nabla_3\cdot\mathbf F_{h,{\rm sub}}
=
-2\overline{\mathbf E'\cdot\mathbf B'}.
\end{equation}
On a reduced brane control volume,
\begin{equation}
\label{eq:part08-stage249-hdot}
\frac{dH_{\rm sub}}{dt}
=
-\Phi_h(t;E)-2\mathcal C_h(t;E),
\end{equation}
where \(\Phi_h\) is the unresolved-helicity flux through the control boundary and \(\mathcal C_h\) is the unresolved covariance source.

The minimal orientation closure introduces \(\sigma=+1\) for aligned and \(\sigma=-1\) for anti-aligned, with
\begin{equation}
\label{eq:part08-stage249-orientation-closure}
\Phi_{h,\sigma}=\Phi_{h,0}+\sigma\Phi_{h,1},
\qquad
\mathcal C_{h,\sigma}=\mathcal C_{h,0}+\sigma\mathcal C_{h,1}.
\end{equation}
Thus
\begin{equation}
\label{eq:part08-stage249-export-rate}
\dot H_{{\rm sub},\sigma}
=\Gamma_0+\sigma\Gamma_1
=\Gamma_0(1+\sigma\alpha_h),
\qquad
\alpha_h=\frac{\Gamma_1}{\Gamma_0}.
\end{equation}
The instantaneous ratio and inverse are
\begin{equation}
\label{eq:part08-stage249-instant-ratio}
R_{\rm inst}=\frac{\dot H_{{\rm sub},+}}{\dot H_{{\rm sub},-}}
=\frac{1+\alpha_h}{1-\alpha_h},
\qquad
\alpha_h=\frac{R_{\rm inst}-1}{R_{\rm inst}+1}.
\end{equation}
Over a common interval,
\begin{equation}
\label{eq:part08-stage249-integrated-ratio}
R_{\rm int}
=\frac{\mathcal H_+}{\mathcal H_-}
=\frac{1+\bar\alpha_h}{1-\bar\alpha_h},
\qquad
\bar\alpha_h=\frac{R_{\rm int}-1}{R_{\rm int}+1}.
\end{equation}
Any common export scale cancels from these ratios.

The benchmark packet is
\begin{equation}
\label{eq:part08-stage249-benchmark-packet}
\begin{aligned}
(\Xi_{\rm turn},\lambda_{\rm th},R_{\rm pk})
&=
(0.34437471,0.42826825,4.94653917),\\
(R_{\rm int},\alpha_{\rm pk},\bar\alpha_h)
&=
(4.10920923,0.66366992,0.60854999).
\end{aligned}
\end{equation}
The status is deliberately limited: this proves a hidden-channel unloading preference inside the declared closure, not intrinsic spin or a Pauli-sector theorem.

\section{Crossing-vs-collapse Goldilocks window}
\label{sec:part08_relaxed_branch:crossing-vs-collapse-goldilocks-window}

Stage 250 turns the dynamic branch into a survival inequality.  The exact path-time integral from the Stage 248 energy law is
\begin{equation}
\label{eq:part08-stage250-path-time}
T_{\rm traj}(E;r_a\to r_b)
=
\sqrt{\frac{m_s}{2}}
\int_{r_b}^{r_a}\frac{dr}{\sqrt{E-V(r)}}.
\end{equation}
The characteristic-width closure used by Session III compresses the active barrier region to
\begin{equation}
\label{eq:part08-stage250-cross-time}
t_{\rm cross}(E)
=
\lambda_{\rm eff}\sqrt{\frac{m_s}{2(E-V_{\rm peak})}}.
\end{equation}
This is a reduced timescale closure, not a replacement for the exact path integral.

The active dressing leg uses the local unstable normal form
\begin{equation}
\label{eq:part08-stage250-v-leg}
\mu_\eta\ddot V=g_{UV}\chi_{\rm peak}V.
\end{equation}
Hence
\begin{equation}
\label{eq:part08-stage250-collapse-time}
t_{\rm collapse}=
\sqrt{\frac{\mu_\eta}{g_{UV}\chi_{\rm peak}}}.
\end{equation}
The survival ratio is
\begin{equation}
\label{eq:part08-stage250-survival-ratio}
\mathcal S(E)=\frac{t_{\rm cross}(E)}{t_{\rm collapse}}
=
\lambda_{\rm eff}
\sqrt{\frac{m_sg_{UV}\chi_{\rm peak}}{2\mu_\eta(E-V_{\rm peak})}}.
\end{equation}
The lower safe edge is \cref{eq:part08-main-goldilocks-edge}, and the speed-space identity is \cref{eq:part08-main-goldilocks-speed}.

If \(\mu_\eta=\alpha m_s\), the edge becomes
\begin{equation}
\label{eq:part08-stage250-heavy-edge}
E_{\rm edge}=V_{\rm peak}+\frac{g_{UV}\chi_{\rm peak}\lambda_{\rm eff}^{2}}{2\alpha},
\end{equation}
so simultaneous scaling of particle mass and wall inertia cancels out of the edge.  The real control packet is
\begin{equation}
\label{eq:part08-stage250-control-packet}
(\lambda_{\rm eff},\chi_{\rm peak},g_{UV},\mu_\eta/m_s).
\end{equation}
With the trigger-width specialization,
\begin{equation}
\label{eq:part08-stage250-trigger-width}
\lambda_{\rm eff}=\lambda_{\rm th}(r_{\rm turn})=
\left|\frac{V(r_{\rm turn})}{V'(r_{\rm turn})}\right|.
\end{equation}
The benchmark edge and sensitivity statements are therefore direct consequences of the dependence
\begin{equation}
\label{eq:part08-stage250-edge-sensitivity}
E_{\rm edge}-V_{\rm peak}\propto\lambda_{\rm eff}^{2}\chi_{\rm peak}.
\end{equation}

Within this closure the safe window is one-sided because \(t_{\rm collapse}\) is energy-independent and \(t_{\rm cross}(E)\) decreases monotonically.  Any upper edge requires an additional high-energy failure mechanism outside Stage 250.

\section{Microscopic export kernel}
\label{sec:part08_relaxed_branch:microscopic-export-kernel}

Stage 251 replaces the phenomenological \(\gamma_{\rm tot}\dot V\) envelope law by the first microscopic odd export operator seen by the active dressing leg.  Project
\begin{equation}
\label{eq:part08-stage251-v-projection}
V(t)=\Pi_{V0}q_0(t)+\Pi_{V-}q_-(t)+\cdots,
\end{equation}
where \(q_0\) is the derivative-coupled scalar outlet and \(q_-\) is the selected mixed quadrupole outlet.

For the scalar lane, take
\begin{equation}
\label{eq:part08-stage251-scalar-coupling}
g_{W,0}(\omega)=\eta_0\omega,
\qquad
g_{A,0}(\omega)=0,
\qquad
\Pi_0^{\rm out}(\omega)=i\gamma_1\omega+O(\omega^3).
\end{equation}
With
\begin{equation}
\label{eq:part08-stage251-delta0}
\Delta_0=\Omega_{U,0}^2\Omega_{W,0}^2-R_0^2,
\end{equation}
the scalar outgoing factor gives
\begin{equation}
\label{eq:part08-stage251-gamma30}
\Gamma_{3,0}=\gamma_1\eta_0^2\frac{\Omega_{U,0}^4}{\Delta_0^2},
\qquad
\Gamma_3=\Pi_{V0}^{2}\Gamma_{3,0}.
\end{equation}
For the selected quadrupole outlet,
\begin{equation}
\label{eq:part08-stage251-gamma5}
\Gamma_{5,-}=\frac{a^5}{27c_s^5}P_{0,-},
\qquad
P_{0,-}=\frac{\beta_0s_-}{\lambda_-},
\qquad
\Gamma_5=\Pi_{V-}^{2}\Gamma_{5,-}.
\end{equation}
Thus the microscopic kernel is \cref{eq:part08-main-kexp}.

The active-leg equation is
\begin{equation}
\label{eq:part08-stage251-active-equation}
\mu_\eta\ddot V-\kappa_VV+
\Gamma_3V^{(3)}-
\Gamma_5V^{(5)}=
\mathcal S_V(t),
\qquad
\kappa_V=g_{UV}\chi_{\rm peak}.
\end{equation}
In Laplace form,
\begin{equation}
\label{eq:part08-stage251-laplace-d}
\mathcal D_V(s)=\mu_\eta s^2-\kappa_V+\Gamma_3s^3+\Gamma_5s^5.
\end{equation}
The passive power identity is
\begin{equation}
\label{eq:part08-stage251-power-identity}
\dot V(\Gamma_3V^{(3)}-\Gamma_5V^{(5)})
=
\frac{d}{dt}\mathcal S_{\rm odd}
-
\Gamma_3\ddot V^{2}
-
\Gamma_5\dddot V^{2},
\end{equation}
so the exported power is
\begin{equation}
\label{eq:part08-stage251-export-power}
\mathcal P_{\rm exp}=\Gamma_3\ddot V^{2}+\Gamma_5\dddot V^{2}\ge0.
\end{equation}

For homogeneous growth \(V=V_0e^{st}\), the characteristic polynomial is
\begin{equation}
\label{eq:part08-stage251-characteristic}
F(s)=\Gamma_5s^5+\Gamma_3s^3+\mu_\eta s^2-\kappa_V.
\end{equation}
For positive \((\Gamma_3,\Gamma_5,\mu_\eta,\kappa_V)\), \(F(0)<0\), \(F(s)\to+\infty\), and \(F'(s)>0\) for \(s>0\).  Therefore there is exactly one positive growth root.  The small-kernel slowdown is
\begin{equation}
\label{eq:part08-stage251-small-kernel}
s_+=s_0-\frac{\Gamma_3s_0^2+\Gamma_5s_0^4}{2\mu_\eta}+O(\Gamma^2),
\qquad
s_0=\sqrt{\kappa_V/\mu_\eta}.
\end{equation}
Thus finite passive cubic/quintic export slows but does not remove the unstable root in the minimal closure.

For an event with crossing rate \(s_c=1/t_{\rm cross}\), strict monotonicity of \(F\) gives the safe half-plane \cref{eq:part08-main-safe-half-plane}.  In the benchmark cold event, this becomes
\begin{equation}
\label{eq:part08-stage251-benchmark-half-plane}
\widehat\Gamma_3+0.3013336471\widehat\Gamma_5\ge 289.61004918.
\end{equation}
This is the exact event-level replacement for the old \(\gamma_{\rm safe}\) inequality.

\section{Vacuum/lattice heat partition and material thresholds}
\label{sec:part08_relaxed_branch:vacuum-lattice-heat-partition-and-material-thresho}

Stages 252 and 253 form the physical companion to the export kernel.  Stage 252 keeps the reduced microscopic language; Stage 253 maps it into physical threshold variables.

\subsection{Channel-resolved heat partition}
\label{subsec:part08-stage252-heat-partition}

Split the microscopic coefficients as
\begin{equation}
\label{eq:part08-stage252-split}
\Gamma_3=\Gamma_3^{\rm vac}+\Gamma_3^{\rm lat},
\qquad
\Gamma_5=\Gamma_5^{\rm vac}+\Gamma_5^{\rm lat},
\qquad
\Gamma_n^{\rm vac},\Gamma_n^{\rm lat}\ge0.
\end{equation}
For an event window \([0,T]\), define
\begin{equation}
\label{eq:part08-stage252-shape-integrals}
\mathcal I_2(T)=\int_0^T\ddot V^{2}\,dt,
\qquad
\mathcal I_3(T)=\int_0^T\dddot V^{2}\,dt.
\end{equation}
The integrated energies are
\begin{equation}
\label{eq:part08-stage252-evac-elat}
E_{\rm vac}=\Gamma_3^{\rm vac}\mathcal I_2+\Gamma_5^{\rm vac}\mathcal I_3,
\qquad
E_{\rm lat}=\Gamma_3^{\rm lat}\mathcal I_2+\Gamma_5^{\rm lat}\mathcal I_3.
\end{equation}
With \(\mathfrak r_V=\mathcal I_3/\mathcal I_2\), the fractions are
\begin{equation}
\label{eq:part08-stage252-fractions}
f_{\rm vac}(\mathfrak r_V)=
\frac{\Gamma_3^{\rm vac}+\Gamma_5^{\rm vac}\mathfrak r_V}{\Gamma_3+\Gamma_5\mathfrak r_V},
\qquad
f_{\rm lat}(\mathfrak r_V)=
\frac{\Gamma_3^{\rm lat}+\Gamma_5^{\rm lat}\mathfrak r_V}{\Gamma_3+\Gamma_5\mathfrak r_V}.
\end{equation}
The partition drift obeys
\begin{equation}
\label{eq:part08-stage252-fraction-drift}
\frac{df_{\rm lat}}{d\mathfrak r_V}
=
\frac{\Gamma_5^{\rm lat}\Gamma_3^{\rm vac}-\Gamma_3^{\rm lat}\Gamma_5^{\rm vac}}{(\Gamma_3+\Gamma_5\mathfrak r_V)^2}.
\end{equation}
So slow events sample the cubic split, fast rough events sample the quintic split.

For a single-growth cold event
\begin{equation}
\label{eq:part08-stage252-single-growth}
V(t)=V_{\rm in}e^{st},
\qquad
0\le t\le T,
\end{equation}
\begin{equation}
\label{eq:part08-stage252-single-growth-integrals}
\mathcal I_2=\frac{V_{\rm in}^{2}s^3}{2}(e^{2sT}-1),
\qquad
\mathcal I_3=s^2\mathcal I_2,
\qquad
\mathfrak r_V=s^2.
\end{equation}
The event-equivalent rates are
\begin{equation}
\label{eq:part08-stage252-equivalent-rates}
\gamma_{\rm vac}^{\rm eq}(s)=\Gamma_3^{\rm vac}s^2+\Gamma_5^{\rm vac}s^4,
\qquad
\gamma_{\rm lat}^{\rm eq}(s)=\Gamma_3^{\rm lat}s^2+\Gamma_5^{\rm lat}s^4.
\end{equation}
At the safe edge, with \(T=1/s_c\) and \(\Gamma_3s_c^3+\Gamma_5s_c^5=\mu_\eta(s_0^2-s_c^2)\), the total minimum exported energy is \cref{eq:part08-main-safe-energy}, and the channel energies are obtained by multiplying by \(f_{\rm vac}(s_c)\) or \(f_{\rm lat}(s_c)\).

The old Session-IV \(3:1\) heat split becomes a coefficient-surface condition:
\begin{equation}
\label{eq:part08-stage252-three-one-surface}
f_{\rm lat}(s_c)=\frac34
\iff
\Gamma_3^{\rm lat}+\Gamma_5^{\rm lat}s_c^2
=
3\bigl(\Gamma_3^{\rm vac}+\Gamma_5^{\rm vac}s_c^2\bigr).
\end{equation}
It is speed-independent only on the special family where both cubic and quintic channels have the same lattice fraction.

\subsection{Physical calibration dictionary}
\label{subsec:part08-stage253-calibration}

Stage 253 introduces the explicit physical dictionary
\begin{equation}
\label{eq:part08-stage253-physical-dictionary}
t^{\rm phys}=t_*t,
\qquad
r^{\rm phys}=\frac{\lambda_{\rm phys}}{\lambda_{\rm ref}}r,
\qquad
E^{\rm phys}=E_*E.
\end{equation}
The safe-edge lattice event-equivalent rate is
\begin{equation}
\label{eq:part08-stage253-lat-safe-rate}
\gamma_{\rm lat,safe}^{\rm eq}
=
f_{\rm lat}(s_c)\mu_\eta\frac{s_0^2-s_c^2}{s_c}.
\end{equation}
A coarse transport projection factor \(\Upsilon_{\rm lat}>0\) defines the physical lattice-turnover rate
\begin{equation}
\label{eq:part08-stage253-lattice-turn-rate}
\gamma_{\rm lat,turn}^{\rm phys}
=
\frac{\gamma_{\rm lat,safe}^{\rm eq}}{\Upsilon_{\rm lat}t_*}
=
\zeta_{\rm ep}\lambda_{\rm ep}\omega_D.
\end{equation}
Thus the electron-phonon turnover threshold is \cref{eq:part08-main-ep-threshold}, and the product form is
\begin{equation}
\label{eq:part08-stage253-ep-product}
(\lambda_{\rm ep}\omega_Dt_{\rm cross}^{\rm phys})_{\min}^{(253)}
=
\frac{f_{\rm lat}(s_c)\mu_\eta(s_0^2-s_c^2)}{\Upsilon_{\rm lat}\zeta_{\rm ep}s_c^2}.
\end{equation}
The legacy Session-V slice is recovered by choosing
\begin{equation}
\label{eq:part08-stage253-upsilon-session}
\Upsilon_{\rm lat}^{\rm sess}
=
\frac{\gamma_{\rm lat,safe}^{\rm eq}}{\gamma_{\rm lattice}^{\rm red}}.
\end{equation}
This is a calibration slice, not a new branch law.

For the harmonic interstitial companion,
\begin{equation}
\label{eq:part08-stage253-physical-radius}
r_{\rm turn}^{\rm phys}=\frac{r_{\rm turn}}{\lambda_{\rm ref}}\lambda_{\rm phys},
\end{equation}
and
\begin{equation}
\label{eq:part08-stage253-chi-lattice}
\chi_{\lambda,{\rm lattice}}(r_{\rm turn})
=
\frac{2\lambda_{\rm phys}}{r_{\rm turn}^{\rm phys}}
=
\frac{2\lambda_{\rm ref}}{r_{\rm turn}}.
\end{equation}
This is only a geometric trigger.  The force-matched stiffness requires
\begin{equation}
\label{eq:part08-stage253-force-match}
k_{\rm eff,req}
=
\frac{E_*\lambda_{\rm ref}^{2}}{\lambda_{\rm phys}^{2}}
\frac{|V_{\rm red}'(r_{\rm turn})|}{r_{\rm turn}}
=
\mathcal K_{\rm turn}\frac{E_*}{\lambda_{\rm phys}^{2}},
\end{equation}
where
\begin{equation}
\label{eq:part08-stage253-kturn}
\mathcal K_{\rm turn}=\frac{\lambda_{\rm ref}^{2}}{r_{\rm turn}}|V_{\rm red}'(r_{\rm turn})|.
\end{equation}
For \(\lambda_{\rm phys}=a_{\rm int}/2\),
\begin{equation}
\label{eq:part08-stage253-k-aint}
k_{\rm eff,req}=4\mathcal K_{\rm turn}\frac{E_*}{a_{\rm int}^{2}}.
\end{equation}

The Korringa companion is the cleanest physical map.  With
\begin{equation}
\label{eq:part08-stage253-korringa}
T_1T=\mathcal K_{\rm corr},
\qquad
T_1\ge t_{\rm cross}^{\rm phys},
\end{equation}
the ceiling temperature is
\begin{equation}
\label{eq:part08-stage253-tmax}
T_{\max}=\frac{\mathcal K_{\rm corr}}{t_{\rm cross}^{\rm phys}}
=\frac{s_c\mathcal K_{\rm corr}}{t_*}.
\end{equation}

\subsection{Material-screening ratios}
\label{subsec:part08-stage253-screening-ratios}

The companion ends with four screening ratios:
\begin{equation}
\label{eq:part08-stage253-pi-ep}
\Pi_{\rm ep}
=
\frac{\Upsilon_{\rm lat}\zeta_{\rm ep}\lambda_{\rm ep}\omega_Dt_*}{\gamma_{\rm lat,safe}^{\rm eq}}
=
\frac{\lambda_{\rm ep}\omega_D}{(\lambda_{\rm ep}\omega_D)_{\min}^{(253)}},
\end{equation}
\begin{equation}
\label{eq:part08-stage253-pi-chi}
\Pi_\chi=\frac{2\lambda_{\rm ref}}{r_{\rm turn}},
\end{equation}
\begin{equation}
\label{eq:part08-stage253-pi-k}
\Pi_k=
\frac{k_{\rm eff}\lambda_{\rm phys}^{2}}{\mathcal K_{\rm turn}E_*}
=
\frac{k_{\rm eff}}{k_{\rm eff,req}},
\end{equation}
\begin{equation}
\label{eq:part08-stage253-pi-t}
\Pi_T(T)=
\frac{\mathcal K_{\rm corr}}{Tt_{\rm cross}^{\rm phys}}
=
\frac{T_{\max}}{T}.
\end{equation}
The materials companion is then
\begin{equation}
\label{eq:part08-stage253-final-screen}
\Pi_{\rm ep}\ge1,
\qquad
\Pi_\chi\ge1,
\qquad
\Pi_k\ge1,
\qquad
\Pi_T(T)\ge1.
\end{equation}
This is not a claim that any particular host passes.  It is the reduced threshold stack that a candidate host must satisfy after the physical unit map and material constants are supplied.

\section{Citation and downstream-use guardrails}
\label{sec:part08-downstream-guardrails}

Part VIII is useful because it lets downstream papers talk about the relaxed branch without contaminating the strict branch.  The safe citation language is:
\begin{quote}
The relaxed/open-system companion is a codimension-three lift of the strict selected-support/orbit front end.  It has an exact recovery map to the strict slice, preserves the no-new-long-range same-charge verdict, and supplies separate leakage/work, non-rigid dressing, compensated-source, barrier, dynamic-event, export-kernel, heat-partition, and material-screening compilers.
\end{quote}

The unsafe citation language is:
\begin{quote}
The relaxed barrier branch proves the strict selected same-charge branch, or supplies a new long-range attraction.
\end{quote}
That statement is false under the status discipline of this ledger.  The relaxed branch is a companion extension around the strict front end.  It is valuable precisely because it states what can change when open-system and non-rigid effects are allowed while preserving what cannot change: the same-charge long-range kernel verdict.

The active open tasks after this part are also explicit.  The completed moving-throat PDE still has to realize or reject the strict selected support/orbit/outgoing packet from Part VII.  If the relaxed companion is used, the actual branch must additionally return the leakage amplitude, non-rigid \((U,V)\) packet, compensated source packet, microscopic export coefficients \((\Gamma_3,\Gamma_5)\) and their vacuum/lattice split, the unit map \((t_*,\lambda_{\rm phys},E_*)\), and the coarse transport factor \(\Upsilon_{\rm lat}\).  Without those realized data, Part VIII remains an exact compiler rather than a completed physical-branch theorem.
